% Options for packages loaded elsewhere
\PassOptionsToPackage{unicode}{hyperref}
\PassOptionsToPackage{hyphens}{url}
\PassOptionsToPackage{dvipsnames,svgnames,x11names}{xcolor}
\documentclass[
  brazilian,
  a4paper,12pt,oneside]{report}
\usepackage{xcolor}
\usepackage[margin=2.5cm]{geometry}
\usepackage{amsmath,amssymb}
\setcounter{secnumdepth}{-\maxdimen} % remove section numbering
\usepackage{iftex}
\ifPDFTeX
  \usepackage[T1]{fontenc}
  \usepackage[utf8]{inputenc}
  \usepackage{textcomp} % provide euro and other symbols
\else % if luatex or xetex
  \usepackage{unicode-math} % this also loads fontspec
  \defaultfontfeatures{Scale=MatchLowercase}
  \defaultfontfeatures[\rmfamily]{Ligatures=TeX,Scale=1}
\fi
\usepackage{lmodern}
\ifPDFTeX\else
  % xetex/luatex font selection
\fi
% Use upquote if available, for straight quotes in verbatim environments
\IfFileExists{upquote.sty}{\usepackage{upquote}}{}
\IfFileExists{microtype.sty}{% use microtype if available
  \usepackage[]{microtype}
  \UseMicrotypeSet[protrusion]{basicmath} % disable protrusion for tt fonts
}{}
\makeatletter
\@ifundefined{KOMAClassName}{% if non-KOMA class
  \IfFileExists{parskip.sty}{%
    \usepackage{parskip}
  }{% else
    \setlength{\parindent}{0pt}
    \setlength{\parskip}{6pt plus 2pt minus 1pt}}
}{% if KOMA class
  \KOMAoptions{parskip=half}}
\makeatother
\usepackage{color}
\usepackage{fancyvrb}
\newcommand{\VerbBar}{|}
\newcommand{\VERB}{\Verb[commandchars=\\\{\}]}
\DefineVerbatimEnvironment{Highlighting}{Verbatim}{commandchars=\\\{\}}
% Add ',fontsize=\small' for more characters per line
\newenvironment{Shaded}{}{}
\newcommand{\AlertTok}[1]{\textcolor[rgb]{1.00,0.00,0.00}{\textbf{#1}}}
\newcommand{\AnnotationTok}[1]{\textcolor[rgb]{0.38,0.63,0.69}{\textbf{\textit{#1}}}}
\newcommand{\AttributeTok}[1]{\textcolor[rgb]{0.49,0.56,0.16}{#1}}
\newcommand{\BaseNTok}[1]{\textcolor[rgb]{0.25,0.63,0.44}{#1}}
\newcommand{\BuiltInTok}[1]{\textcolor[rgb]{0.00,0.50,0.00}{#1}}
\newcommand{\CharTok}[1]{\textcolor[rgb]{0.25,0.44,0.63}{#1}}
\newcommand{\CommentTok}[1]{\textcolor[rgb]{0.38,0.63,0.69}{\textit{#1}}}
\newcommand{\CommentVarTok}[1]{\textcolor[rgb]{0.38,0.63,0.69}{\textbf{\textit{#1}}}}
\newcommand{\ConstantTok}[1]{\textcolor[rgb]{0.53,0.00,0.00}{#1}}
\newcommand{\ControlFlowTok}[1]{\textcolor[rgb]{0.00,0.44,0.13}{\textbf{#1}}}
\newcommand{\DataTypeTok}[1]{\textcolor[rgb]{0.56,0.13,0.00}{#1}}
\newcommand{\DecValTok}[1]{\textcolor[rgb]{0.25,0.63,0.44}{#1}}
\newcommand{\DocumentationTok}[1]{\textcolor[rgb]{0.73,0.13,0.13}{\textit{#1}}}
\newcommand{\ErrorTok}[1]{\textcolor[rgb]{1.00,0.00,0.00}{\textbf{#1}}}
\newcommand{\ExtensionTok}[1]{#1}
\newcommand{\FloatTok}[1]{\textcolor[rgb]{0.25,0.63,0.44}{#1}}
\newcommand{\FunctionTok}[1]{\textcolor[rgb]{0.02,0.16,0.49}{#1}}
\newcommand{\ImportTok}[1]{\textcolor[rgb]{0.00,0.50,0.00}{\textbf{#1}}}
\newcommand{\InformationTok}[1]{\textcolor[rgb]{0.38,0.63,0.69}{\textbf{\textit{#1}}}}
\newcommand{\KeywordTok}[1]{\textcolor[rgb]{0.00,0.44,0.13}{\textbf{#1}}}
\newcommand{\NormalTok}[1]{#1}
\newcommand{\OperatorTok}[1]{\textcolor[rgb]{0.40,0.40,0.40}{#1}}
\newcommand{\OtherTok}[1]{\textcolor[rgb]{0.00,0.44,0.13}{#1}}
\newcommand{\PreprocessorTok}[1]{\textcolor[rgb]{0.74,0.48,0.00}{#1}}
\newcommand{\RegionMarkerTok}[1]{#1}
\newcommand{\SpecialCharTok}[1]{\textcolor[rgb]{0.25,0.44,0.63}{#1}}
\newcommand{\SpecialStringTok}[1]{\textcolor[rgb]{0.73,0.40,0.53}{#1}}
\newcommand{\StringTok}[1]{\textcolor[rgb]{0.25,0.44,0.63}{#1}}
\newcommand{\VariableTok}[1]{\textcolor[rgb]{0.10,0.09,0.49}{#1}}
\newcommand{\VerbatimStringTok}[1]{\textcolor[rgb]{0.25,0.44,0.63}{#1}}
\newcommand{\WarningTok}[1]{\textcolor[rgb]{0.38,0.63,0.69}{\textbf{\textit{#1}}}}
\usepackage{longtable,booktabs,array}
\newcounter{none} % for unnumbered tables
\usepackage{calc} % for calculating minipage widths
% Correct order of tables after \paragraph or \subparagraph
\usepackage{etoolbox}
\makeatletter
\patchcmd\longtable{\par}{\if@noskipsec\mbox{}\fi\par}{}{}
\makeatother
% Allow footnotes in longtable head/foot
\IfFileExists{footnotehyper.sty}{\usepackage{footnotehyper}}{\usepackage{footnote}}
\makesavenoteenv{longtable}
\ifLuaTeX
\usepackage[bidi=basic,shorthands=off]{babel}
\else
\usepackage[bidi=default,shorthands=off]{babel}
\fi
\ifLuaTeX
  \usepackage{selnolig} % disable illegal ligatures
\fi
\setlength{\emergencystretch}{3em} % prevent overfull lines
\providecommand{\tightlist}{%
  \setlength{\itemsep}{0pt}\setlength{\parskip}{0pt}}
\usepackage{bookmark}
\IfFileExists{xurl.sty}{\usepackage{xurl}}{} % add URL line breaks if available
\urlstyle{same}
\hypersetup{
  pdftitle={Manual de Configuração},
  pdfauthor={Documentação do sistema AGIA},
  pdflang={pt-BR},
  colorlinks=true,
  linkcolor={blue},
  filecolor={Maroon},
  citecolor={Blue},
  urlcolor={blue},
  pdfcreator={LaTeX via pandoc}}

\title{Manual de Configuração}
\usepackage{etoolbox}
\makeatletter
\providecommand{\subtitle}[1]{% add subtitle to \maketitle
  \apptocmd{\@title}{\par {\large #1 \par}}{}{}
}
\makeatother
\subtitle{Sistema AGIA --- Prospecção Tecnológica}
\author{Documentação do sistema AGIA}
\date{17/10/2026}

\begin{document}
\maketitle

{
\setcounter{tocdepth}{1}
\tableofcontents
}
\emph{Atualizado em 17/10/2026.}

Este manual descreve \textbf{tudo o que pode ser configurado} no
sistema: variáveis do arquivo \texttt{.}\allowbreak{}\texttt{env},
parâmetros da tela \textbf{Configurações}, provedores de IA, APIs de
busca, relatório e segurança.

Pré-requisito: o sistema instalado e em execução, com o banco migrado e
o backend no ar.

\begin{center}\rule{0.5\linewidth}{0.5pt}\end{center}

\chapter{1. Como a configuração é
organizada}\label{como-a-configurauxe7uxe3o-uxe9-organizada}

Há \textbf{três lugares} onde um valor pode ser definido. Eles não são
equivalentes:

{\def\LTcaptype{none} % do not increment counter
\begin{longtable}[]{@{}
  >{\raggedright\arraybackslash}p{(\linewidth - 4\tabcolsep) * \real{0.3333}}
  >{\raggedright\arraybackslash}p{(\linewidth - 4\tabcolsep) * \real{0.3333}}
  >{\raggedright\arraybackslash}p{(\linewidth - 4\tabcolsep) * \real{0.3333}}@{}}
\toprule\noalign{}
\begin{minipage}[b]{\linewidth}\raggedright
Onde
\end{minipage} & \begin{minipage}[b]{\linewidth}\raggedright
O que fica ali
\end{minipage} & \begin{minipage}[b]{\linewidth}\raggedright
Quando vale
\end{minipage} \\
\midrule\noalign{}
\endhead
\bottomrule\noalign{}
\endlastfoot
\textbf{Arquivo \texttt{.}\allowbreak{}\texttt{env}} (na raiz do
projeto) & Infraestrutura (banco, MinIO, endereços dos serviços), flags
de API e \textbf{credenciais} & Infraestrutura e flags: lidas sempre.
Credenciais: copiadas para o banco só na \textbf{primeira subida}. Não
definem mais os parâmetros de comportamento. \\
\textbf{Padrões do código} & Os valores iniciais de todos os parâmetros
de comportamento & Gravados no banco pelo seed automático, na primeira
subida. \\
\textbf{Banco de dados} (tabelas
\texttt{app\_}\allowbreak{}\texttt{settings},
\texttt{search\_}\allowbreak{}\texttt{api\_}\allowbreak{}\texttt{selection},
\texttt{llm\_}\allowbreak{}\texttt{provider\_}\allowbreak{}\texttt{configs},
\texttt{llm\_}\allowbreak{}\texttt{call\_}\allowbreak{}\texttt{site\_}\allowbreak{}\texttt{bindings})
& Todas as configurações editáveis & Fonte da verdade depois do seed. \\
\textbf{Tela Configurações} (ou API
\texttt{/}\allowbreak{}\texttt{config/}\allowbreak{}\texttt{*}) &
Interface para editar o banco & Onde você deve mudar valores no dia a
dia. \\
\end{longtable}
}

\textbf{Regra prática:} para mudar um parâmetro de comportamento ou um
modelo de IA, use a \textbf{tela Configurações}. O
\texttt{.}\allowbreak{}\texttt{env} não controla mais esses valores, em
nenhuma subida. Flags de API continuam lidas do
\texttt{.}\allowbreak{}\texttt{env} a cada requisição. Credenciais são
copiadas do \texttt{.}\allowbreak{}\texttt{env} para o banco na primeira
subida; depois disso, altere-as pela tela.

\begin{center}\rule{0.5\linewidth}{0.5pt}\end{center}

\chapter{\texorpdfstring{2. Ciclo de vida: do
\texttt{.}\allowbreak{}\texttt{env} ao
banco}{2. Ciclo de vida: do .env ao banco}}\label{ciclo-de-vida-do-.env-ao-banco}

Na subida do backend:

\begin{enumerate}
\def\labelenumi{\arabic{enumi}.}
\tightlist
\item
  \textbf{Migrações} (\texttt{alembic\ upgrade\ head}, passo manual da
  instalação) criam as tabelas de sessão.
\item
  O backend cria as demais tabelas e executa o \textbf{seed}, que
  preenche as tabelas de configuração.
\item
  O \textbf{seed} só grava valores de cada tabela de configuração
  \textbf{se ela estiver vazia}. Os parâmetros de comportamento vêm dos
  padrões do código; credenciais vêm do
  \texttt{.}\allowbreak{}\texttt{env}. Chaves novas introduzidas em
  versões futuras são acrescentadas sem sobrescrever o que já existe.
\item
  \textbf{Sincronização:} os valores do banco são carregados na memória
  do backend, e o código passa a usá-los a partir daí.
\item
  \textbf{Edição pela tela:} cada alteração é gravada no banco e
  aplicada \textbf{imediatamente}, sem reiniciar, exceto nos casos
  listados na \hyperref[10-quando-reiniciar-o-backend]{seção 10}.
\end{enumerate}

\textbf{Consequência importante:} editar parâmetros de comportamento no
\texttt{.}\allowbreak{}\texttt{env} não tem efeito. Para mudar um valor,
use a tela. Para voltar ao padrão do código, apague a linha da chave na
tabela \texttt{app\_}\allowbreak{}\texttt{settings} e reinicie o
backend: o seed recria a chave com o padrão.

\begin{quote}
Se o seu banco foi criado antes desta mudança, alguns valores podem ter
vindo do \texttt{.}\allowbreak{}\texttt{env} daquela época (ex.: ano
inicial 2010, relevância 0,5). Eles continuam no banco e valem até você
alterá-los. Confira na tela \textbf{Configurações} o valor vigente.
\end{quote}

\begin{center}\rule{0.5\linewidth}{0.5pt}\end{center}

\chapter{\texorpdfstring{3. Arquivo
\texttt{.}\allowbreak{}\texttt{env}}{3. Arquivo .env}}\label{arquivo-.env}

O \texttt{.}\allowbreak{}\texttt{env} é lido pelo sistema na subida e
\textbf{não deve ser versionado} (nem enviado a repositórios), porque
contém credenciais. Ele é criado na instalação, com as variáveis
descritas abaixo.

Convenção: o nome da variável é o nome do parâmetro em
\texttt{UPPER\_}\allowbreak{}\texttt{SNAKE\_}\allowbreak{}\texttt{CASE}.
Ex.:
\texttt{ollama\_}\allowbreak{}\texttt{base\_}\allowbreak{}\texttt{url}
-\textgreater{}
\texttt{OLLAMA\_}\allowbreak{}\texttt{BASE\_}\allowbreak{}\texttt{URL}.

\begin{quote}
\textbf{Só estas variáveis são lidas do
\texttt{.}\allowbreak{}\texttt{env}:} infraestrutura (banco, MinIO,
ChromaDB, LaTeX, LanguageTool, endereço do Ollama), flags de API e
credenciais. Credenciais são copiadas para o banco só na primeira
subida. Os parâmetros de comportamento (busca, IA, extração de termos,
inferência, relatório) não estão aqui: seus padrões estão na
\hyperref[6-paruxe2metros-da-tela-configurauxe7uxf5es]{seção 6} e são
alterados pela tela.
\end{quote}

\section{3.1 Banco de dados}\label{banco-de-dados}

{\def\LTcaptype{none} % do not increment counter
\begin{longtable}[]{@{}
  >{\raggedright\arraybackslash}p{(\linewidth - 6\tabcolsep) * \real{0.0833}}
  >{\raggedright\arraybackslash}p{(\linewidth - 6\tabcolsep) * \real{0.3333}}
  >{\raggedright\arraybackslash}p{(\linewidth - 6\tabcolsep) * \real{0.3333}}
  >{\raggedright\arraybackslash}p{(\linewidth - 6\tabcolsep) * \real{0.2500}}@{}}
\toprule\noalign{}
\begin{minipage}[b]{\linewidth}\raggedright
Variável
\end{minipage} & \begin{minipage}[b]{\linewidth}\raggedright
Padrão
\end{minipage} & \begin{minipage}[b]{\linewidth}\raggedright
Descrição
\end{minipage} & \begin{minipage}[b]{\linewidth}\raggedright
Quando mudar
\end{minipage} \\
\midrule\noalign{}
\endhead
\bottomrule\noalign{}
\endlastfoot
\texttt{DATABASE\_}\allowbreak{}\texttt{URL} &
\texttt{postgresql+asyncpg:}\allowbreak{}\texttt{/}\allowbreak{}\texttt{/}\allowbreak{}\texttt{postgres:}\allowbreak{}\texttt{postgres@localhost:}\allowbreak{}\texttt{5432/}\allowbreak{}\texttt{pfc\_}\allowbreak{}\texttt{db}
& URL do PostgreSQL. O driver \texttt{asyncpg} é usado pelo backend; as
migrações usam \texttt{psycopg2} automaticamente. & Banco em outro host,
porta, usuário ou senha. \\
\end{longtable}
}

\section{3.2 Aplicação}\label{aplicauxe7uxe3o}

{\def\LTcaptype{none} % do not increment counter
\begin{longtable}[]{@{}
  >{\raggedright\arraybackslash}p{(\linewidth - 4\tabcolsep) * \real{0.1311}}
  >{\raggedright\arraybackslash}p{(\linewidth - 4\tabcolsep) * \real{0.2787}}
  >{\raggedright\arraybackslash}p{(\linewidth - 4\tabcolsep) * \real{0.5902}}@{}}
\toprule\noalign{}
\begin{minipage}[b]{\linewidth}\raggedright
Variável
\end{minipage} & \begin{minipage}[b]{\linewidth}\raggedright
Padrão
\end{minipage} & \begin{minipage}[b]{\linewidth}\raggedright
Descrição
\end{minipage} \\
\midrule\noalign{}
\endhead
\bottomrule\noalign{}
\endlastfoot
\texttt{APP\_}\allowbreak{}\texttt{NAME} &
\texttt{Technology\ Prospecting\ API} & Nome exibido na documentação da
API e nos logs. \\
\texttt{APP\_}\allowbreak{}\texttt{VERSION} &
\texttt{0.}\allowbreak{}\texttt{1.}\allowbreak{}\texttt{0} & Versão
exibida na documentação da API. \\
\texttt{ENVIRONMENT} & \texttt{development} & Com \texttt{production}, o
banco usa pool de conexões; fora disso, conexões são abertas sob
demanda. Também controla o formato dos logs. \\
\texttt{DEBUG} & \texttt{True} & Modo de depuração do FastAPI. Desligue
em ambientes compartilhados. \\
\texttt{LOG\_}\allowbreak{}\texttt{LEVEL} & \texttt{INFO} & Nível de
log. \\
\texttt{SQL\_}\allowbreak{}\texttt{ECHO} & \texttt{False} & Registra
\textbf{todas} as queries SQL no log. Ligue só para depurar uma query
específica (o log fica muito grande). \\
\end{longtable}
}

\section{3.3 API}\label{api}

{\def\LTcaptype{none} % do not increment counter
\begin{longtable}[]{@{}
  >{\raggedright\arraybackslash}p{(\linewidth - 4\tabcolsep) * \real{0.1186}}
  >{\raggedright\arraybackslash}p{(\linewidth - 4\tabcolsep) * \real{0.4407}}
  >{\raggedright\arraybackslash}p{(\linewidth - 4\tabcolsep) * \real{0.4407}}@{}}
\toprule\noalign{}
\begin{minipage}[b]{\linewidth}\raggedright
Variável
\end{minipage} & \begin{minipage}[b]{\linewidth}\raggedright
Padrão
\end{minipage} & \begin{minipage}[b]{\linewidth}\raggedright
Descrição
\end{minipage} \\
\midrule\noalign{}
\endhead
\bottomrule\noalign{}
\endlastfoot
\texttt{API\_}\allowbreak{}\texttt{PREFIX} &
\texttt{/}\allowbreak{}\texttt{api/}\allowbreak{}\texttt{v1} & Prefixo
de todas as rotas. O frontend assume esse valor. Mudar exige ajustar
\texttt{VITE\_}\allowbreak{}\texttt{API\_}\allowbreak{}\texttt{BASE\_}\allowbreak{}\texttt{URL}. \\
\texttt{ALLOWED\_}\allowbreak{}\texttt{ORIGINS} &
\texttt{{[}"http:}\allowbreak{}\texttt{/}\allowbreak{}\texttt{/}\allowbreak{}\texttt{localhost:}\allowbreak{}\texttt{3000","http:}\allowbreak{}\texttt{/}\allowbreak{}\texttt{/}\allowbreak{}\texttt{localhost:}\allowbreak{}\texttt{5173","http:}\allowbreak{}\texttt{/}\allowbreak{}\texttt{/}\allowbreak{}\texttt{localhost:}\allowbreak{}\texttt{8000"{]}}
& Origens autorizadas pelo CORS. \textbf{Formato JSON obrigatório.}
Inclua a origem do frontend (ex.:
\texttt{http:}\allowbreak{}\texttt{/}\allowbreak{}\texttt{/}\allowbreak{}\texttt{localhost:}\allowbreak{}\texttt{5173}). \\
\end{longtable}
}

\section{3.4 Flags de API}\label{flags-de-api}

{\def\LTcaptype{none} % do not increment counter
\begin{longtable}[]{@{}
  >{\raggedright\arraybackslash}p{(\linewidth - 4\tabcolsep) * \real{0.2667}}
  >{\raggedright\arraybackslash}p{(\linewidth - 4\tabcolsep) * \real{0.0667}}
  >{\raggedright\arraybackslash}p{(\linewidth - 4\tabcolsep) * \real{0.6667}}@{}}
\toprule\noalign{}
\begin{minipage}[b]{\linewidth}\raggedright
Variável
\end{minipage} & \begin{minipage}[b]{\linewidth}\raggedright
Padrão
\end{minipage} & \begin{minipage}[b]{\linewidth}\raggedright
Descrição
\end{minipage} \\
\midrule\noalign{}
\endhead
\bottomrule\noalign{}
\endlastfoot
\texttt{LENS\_}\allowbreak{}\texttt{PATENT\_}\allowbreak{}\texttt{ENABLED}
& \texttt{True} & Habilita a Lens como fonte de \textbf{patentes} na
busca final. \\
\texttt{LENS\_}\allowbreak{}\texttt{SCHOLARLY\_}\allowbreak{}\texttt{ENABLED}
& \texttt{True} & Habilita a Lens como fonte de \textbf{artigos} na
busca final. \\
\texttt{OPS\_}\allowbreak{}\texttt{ENABLED} & \texttt{True} & Habilita a
OPS (Escritório Europeu de Patentes) para patentes. \\
\texttt{SCOPUS\_}\allowbreak{}\texttt{ENABLED} & \texttt{True} &
Habilita a Scopus para artigos. \\
\texttt{LENS\_}\allowbreak{}\texttt{ENABLED} & \texttt{True} &
Liga/desliga \textbf{as duas} APIs Lens de uma vez (compatibilidade com
versões anteriores). \\
\end{longtable}
}

Essas flags são lidas a cada requisição, na montagem das listas de
campos disponíveis para a IA. Na \textbf{primeira subida},
\texttt{OPS\_}\allowbreak{}\texttt{ENABLED} e
\texttt{SCOPUS\_}\allowbreak{}\texttt{ENABLED} também definem qual API
cada família começa ativa. Depois disso, a escolha é feita na tela
Configurações › Busca (ver \hyperref[5-apis-de-busca]{seção 5}).

\section{3.5 MinIO (armazenamento dos gráficos, das figuras e da
capa)}\label{minio-armazenamento-dos-gruxe1ficos-das-figuras-e-da-capa}

{\def\LTcaptype{none} % do not increment counter
\begin{longtable}[]{@{}
  >{\raggedright\arraybackslash}p{(\linewidth - 4\tabcolsep) * \real{0.1864}}
  >{\raggedright\arraybackslash}p{(\linewidth - 4\tabcolsep) * \real{0.1695}}
  >{\raggedright\arraybackslash}p{(\linewidth - 4\tabcolsep) * \real{0.6441}}@{}}
\toprule\noalign{}
\begin{minipage}[b]{\linewidth}\raggedright
Variável
\end{minipage} & \begin{minipage}[b]{\linewidth}\raggedright
Padrão
\end{minipage} & \begin{minipage}[b]{\linewidth}\raggedright
Descrição
\end{minipage} \\
\midrule\noalign{}
\endhead
\bottomrule\noalign{}
\endlastfoot
\texttt{MINIO\_}\allowbreak{}\texttt{ENDPOINT} &
\texttt{localhost:}\allowbreak{}\texttt{9000} & Endereço da API S3 (sem
\texttt{http:}\allowbreak{}\texttt{/}\allowbreak{}\texttt{/}). \\
\texttt{MINIO\_}\allowbreak{}\texttt{ACCESS\_}\allowbreak{}\texttt{KEY}
& \texttt{minioadmin} & Usuário. Deve coincidir com o usuário
configurado no serviço MinIO. \\
\texttt{MINIO\_}\allowbreak{}\texttt{SECRET\_}\allowbreak{}\texttt{KEY}
& \texttt{minioadmin} & Senha. Deve coincidir com a senha configurada no
serviço MinIO. \\
\texttt{MINIO\_}\allowbreak{}\texttt{BUCKET} &
\texttt{session-}\allowbreak{}\texttt{charts} & Bucket criado
automaticamente na subida. \\
\texttt{MINIO\_}\allowbreak{}\texttt{SECURE} & \texttt{False} &
\texttt{True} para HTTPS. \\
\end{longtable}
}

Se o MinIO não responder, o backend sobe com aviso
(\texttt{minio\_}\allowbreak{}\texttt{bucket\_}\allowbreak{}\texttt{ensure\_}\allowbreak{}\texttt{failed}),
mas gráficos, anexos e a capa do relatório não são persistidos.

\section{3.6 ChromaDB (RAG do
relatório)}\label{chromadb-rag-do-relatuxf3rio}

{\def\LTcaptype{none} % do not increment counter
\begin{longtable}[]{@{}lll@{}}
\toprule\noalign{}
Variável & Padrão & Descrição \\
\midrule\noalign{}
\endhead
\bottomrule\noalign{}
\endlastfoot
\texttt{CHROMA\_}\allowbreak{}\texttt{HOST} & \texttt{localhost} &
Endereço do serviço ChromaDB. \\
\texttt{CHROMA\_}\allowbreak{}\texttt{PORT} & \texttt{8001} & Porta no
host. A 8000 é a do backend. \\
\end{longtable}
}

\section{3.7 Compilador LaTeX e revisão de
língua}\label{compilador-latex-e-revisuxe3o-de-luxedngua}

{\def\LTcaptype{none} % do not increment counter
\begin{longtable}[]{@{}
  >{\raggedright\arraybackslash}p{(\linewidth - 4\tabcolsep) * \real{0.2167}}
  >{\raggedright\arraybackslash}p{(\linewidth - 4\tabcolsep) * \real{0.2500}}
  >{\raggedright\arraybackslash}p{(\linewidth - 4\tabcolsep) * \real{0.5333}}@{}}
\toprule\noalign{}
\begin{minipage}[b]{\linewidth}\raggedright
Variável
\end{minipage} & \begin{minipage}[b]{\linewidth}\raggedright
Padrão
\end{minipage} & \begin{minipage}[b]{\linewidth}\raggedright
Descrição
\end{minipage} \\
\midrule\noalign{}
\endhead
\bottomrule\noalign{}
\endlastfoot
\texttt{LATEX\_}\allowbreak{}\texttt{COMPILER\_}\allowbreak{}\texttt{URL}
&
\texttt{http:}\allowbreak{}\texttt{/}\allowbreak{}\texttt{/}\allowbreak{}\texttt{localhost:}\allowbreak{}\texttt{8090}
& Serviço que compila o \texttt{.}\allowbreak{}\texttt{tex} em PDF (sob
demanda). \\
\texttt{LANGUAGETOOL\_}\allowbreak{}\texttt{URL} &
\texttt{http:}\allowbreak{}\texttt{/}\allowbreak{}\texttt{/}\allowbreak{}\texttt{localhost:}\allowbreak{}\texttt{8010}
& Serviço de revisão ortográfica e gramatical. \\
\end{longtable}
}

\section{3.8 Chaves de provedores de IA
externos}\label{chaves-de-provedores-de-ia-externos}

{\def\LTcaptype{none} % do not increment counter
\begin{longtable}[]{@{}
  >{\raggedright\arraybackslash}p{(\linewidth - 4\tabcolsep) * \real{0.2667}}
  >{\raggedright\arraybackslash}p{(\linewidth - 4\tabcolsep) * \real{0.0667}}
  >{\raggedright\arraybackslash}p{(\linewidth - 4\tabcolsep) * \real{0.6667}}@{}}
\toprule\noalign{}
\begin{minipage}[b]{\linewidth}\raggedright
Variável
\end{minipage} & \begin{minipage}[b]{\linewidth}\raggedright
Padrão
\end{minipage} & \begin{minipage}[b]{\linewidth}\raggedright
Descrição
\end{minipage} \\
\midrule\noalign{}
\endhead
\bottomrule\noalign{}
\endlastfoot
\texttt{LLM\_}\allowbreak{}\texttt{GEMINI\_}\allowbreak{}\texttt{API\_}\allowbreak{}\texttt{KEY}
& vazio & Chave do Google Gemini. Copiada para o banco na primeira
subida; depois, edite pela tela. \\
\texttt{LLM\_}\allowbreak{}\texttt{ANTHROPIC\_}\allowbreak{}\texttt{API\_}\allowbreak{}\texttt{KEY}
& vazio & Chave da Anthropic (é separada de uma assinatura Claude Pro).
Mesmo comportamento da chave do Gemini. \\
\end{longtable}
}

O provedor e o modelo de cada ponto de uso são escolhidos na tela (ver
\hyperref[4-provedores-de-ia]{seção 4}). Os modelos iniciais estão na
\hyperref[63-inteliguxeancia-artificial]{seção 6.3}.

\section{3.9 Endereço do Ollama}\label{endereuxe7o-do-ollama}

{\def\LTcaptype{none} % do not increment counter
\begin{longtable}[]{@{}
  >{\raggedright\arraybackslash}p{(\linewidth - 4\tabcolsep) * \real{0.1667}}
  >{\raggedright\arraybackslash}p{(\linewidth - 4\tabcolsep) * \real{0.2333}}
  >{\raggedright\arraybackslash}p{(\linewidth - 4\tabcolsep) * \real{0.6000}}@{}}
\toprule\noalign{}
\begin{minipage}[b]{\linewidth}\raggedright
Variável
\end{minipage} & \begin{minipage}[b]{\linewidth}\raggedright
Padrão
\end{minipage} & \begin{minipage}[b]{\linewidth}\raggedright
Descrição
\end{minipage} \\
\midrule\noalign{}
\endhead
\bottomrule\noalign{}
\endlastfoot
\texttt{OLLAMA\_}\allowbreak{}\texttt{BASE\_}\allowbreak{}\texttt{URL} &
\texttt{http:}\allowbreak{}\texttt{/}\allowbreak{}\texttt{/}\allowbreak{}\texttt{localhost:}\allowbreak{}\texttt{11434}
& Endereço do servidor compatível com a API da OpenAI. Local: o serviço
do Docker. Intranet: o endpoint da intranet. Copiado para o banco na
primeira subida. \\
\texttt{OLLAMA\_}\allowbreak{}\texttt{API\_}\allowbreak{}\texttt{KEY} &
vazio & Bearer token. Deixe vazio para o Ollama local, que não exige
autenticação. \\
\end{longtable}
}

\section{3.10 Credenciais de busca}\label{credenciais-de-busca}

Veja a \hyperref[5-apis-de-busca]{seção 5} para o passo a passo.

{\def\LTcaptype{none} % do not increment counter
\begin{longtable}[]{@{}lll@{}}
\toprule\noalign{}
Variável & Padrão & Descrição \\
\midrule\noalign{}
\endhead
\bottomrule\noalign{}
\endlastfoot
\texttt{LENS\_}\allowbreak{}\texttt{API\_}\allowbreak{}\texttt{TOKEN} &
vazio & Token da API Lens (patentes e artigos). \\
\texttt{OPS\_}\allowbreak{}\texttt{CONSUMER\_}\allowbreak{}\texttt{KEY}
& vazio & Consumer key OAuth2 da OPS (EPO). \\
\texttt{OPS\_}\allowbreak{}\texttt{CONSUMER\_}\allowbreak{}\texttt{SECRET}
& vazio & Consumer secret OAuth2 da OPS. \\
\texttt{SCOPUS\_}\allowbreak{}\texttt{API\_}\allowbreak{}\texttt{KEY} &
vazio & Chave da Scopus (Elsevier). \\
\end{longtable}
}

\section{3.11 Frontend}\label{frontend}

{\def\LTcaptype{none} % do not increment counter
\begin{longtable}[]{@{}
  >{\raggedright\arraybackslash}p{(\linewidth - 4\tabcolsep) * \real{0.1667}}
  >{\raggedright\arraybackslash}p{(\linewidth - 4\tabcolsep) * \real{0.2833}}
  >{\raggedright\arraybackslash}p{(\linewidth - 4\tabcolsep) * \real{0.5500}}@{}}
\toprule\noalign{}
\begin{minipage}[b]{\linewidth}\raggedright
Variável
\end{minipage} & \begin{minipage}[b]{\linewidth}\raggedright
Padrão
\end{minipage} & \begin{minipage}[b]{\linewidth}\raggedright
Descrição
\end{minipage} \\
\midrule\noalign{}
\endhead
\bottomrule\noalign{}
\endlastfoot
\texttt{VITE\_}\allowbreak{}\texttt{API\_}\allowbreak{}\texttt{BASE\_}\allowbreak{}\texttt{URL}
&
\texttt{http:}\allowbreak{}\texttt{/}\allowbreak{}\texttt{/}\allowbreak{}\texttt{localhost:}\allowbreak{}\texttt{8000/}\allowbreak{}\texttt{api/}\allowbreak{}\texttt{v1}
& Endereço da API vista pelo navegador. Só mude se o backend estiver em
outro endereço. Vale após reiniciar o servidor de desenvolvimento do
frontend. \\
\end{longtable}
}

Se essa variável não for definida, o padrão do código é usado.

\begin{center}\rule{0.5\linewidth}{0.5pt}\end{center}

\chapter{4. Provedores de IA}\label{provedores-de-ia}

O sistema usa IA em \textbf{cinco pontos de uso} (chamadas), e cada um
pode apontar para um provedor diferente:

{\def\LTcaptype{none} % do not increment counter
\begin{longtable}[]{@{}
  >{\raggedright\arraybackslash}p{(\linewidth - 2\tabcolsep) * \real{0.2333}}
  >{\raggedright\arraybackslash}p{(\linewidth - 2\tabcolsep) * \real{0.7667}}@{}}
\toprule\noalign{}
\begin{minipage}[b]{\linewidth}\raggedright
Ponto de uso
\end{minipage} & \begin{minipage}[b]{\linewidth}\raggedright
Para que serve
\end{minipage} \\
\midrule\noalign{}
\endhead
\bottomrule\noalign{}
\endlastfoot
\texttt{theme\_}\allowbreak{}\texttt{candidates} & Geração e
especificação de variações do tema (Refinamento de Parâmetros). \\
\texttt{probe\_}\allowbreak{}\texttt{query} & Geração da query
exploratória. \\
\texttt{final\_}\allowbreak{}\texttt{query} & Geração da query final. \\
\texttt{report\_}\allowbreak{}\texttt{writing} & Redação do texto das
seções do relatório. \\
\texttt{report\_}\allowbreak{}\texttt{review} & Sugestões de revisão de
texto do relatório. \\
\end{longtable}
}

\section{4.1 Cadastrar um provedor}\label{cadastrar-um-provedor}

Tela \textbf{Configurações › Inteligência Artificial › Provider}, botão
\textbf{Adicionar nova configuração}. Campos:

\begin{itemize}
\tightlist
\item
  \textbf{Provider:} \texttt{gemini}, \texttt{anthropic} ou
  \texttt{ollama} (este cobre qualquer servidor compatível com a API da
  OpenAI: Ollama local ou o LLM da intranet).
\item
  \textbf{Modelo:} nome do modelo no provedor (ex.:
  \texttt{gemini-}\allowbreak{}\texttt{2.}\allowbreak{}\texttt{0-}\allowbreak{}\texttt{flash-}\allowbreak{}\texttt{exp},
  \texttt{claude-}\allowbreak{}\texttt{haiku-}\allowbreak{}\texttt{4-}\allowbreak{}\texttt{5},
  \texttt{qwen2.}\allowbreak{}\texttt{5:}\allowbreak{}\texttt{3b-}\allowbreak{}\texttt{instruct}).
\item
  \textbf{Base URL:} obrigatória para \texttt{ollama}.
\item
  \textbf{API key:} obrigatória para \texttt{gemini} e
  \texttt{anthropic}; \textbf{opcional} para \texttt{ollama}.
\end{itemize}

Uma mesma combinação de provedor, modelo e URL não pode ser cadastrada
duas vezes (o sistema recusa com ``talvez já exista uma igual'').

\section{4.2 Escolher o provedor de cada ponto de
uso}\label{escolher-o-provedor-de-cada-ponto-de-uso}

Na mesma tela, \textbf{Qual IA cada chamada usa}: para cada ponto de
uso, escolha uma das configurações cadastradas. A troca vale para a
\textbf{próxima} chamada.

\section{4.3 Remover um provedor}\label{remover-um-provedor}

\textbf{Remover configuração} exclui o cadastro. Os pontos de uso que o
usavam precisam ser reapontados antes.

\section{4.4 Chaves de API e
segurança}\label{chaves-de-api-e-seguranuxe7a}

A chave gravada no banco aparece \textbf{mascarada} nas listagens da
tela. Ela fica guardada no banco, então proteja o acesso ao PostgreSQL
(ver \hyperref[9-seguranuxe7a]{seção 9}).

\begin{center}\rule{0.5\linewidth}{0.5pt}\end{center}

\chapter{5. APIs de busca}\label{apis-de-busca}

O sistema consulta \textbf{duas famílias} de fontes:

{\def\LTcaptype{none} % do not increment counter
\begin{longtable}[]{@{}
  >{\raggedright\arraybackslash}p{(\linewidth - 4\tabcolsep) * \real{0.1356}}
  >{\raggedright\arraybackslash}p{(\linewidth - 4\tabcolsep) * \real{0.5763}}
  >{\raggedright\arraybackslash}p{(\linewidth - 4\tabcolsep) * \real{0.2881}}@{}}
\toprule\noalign{}
\begin{minipage}[b]{\linewidth}\raggedright
Família
\end{minipage} & \begin{minipage}[b]{\linewidth}\raggedright
Fontes disponíveis
\end{minipage} & \begin{minipage}[b]{\linewidth}\raggedright
Padrão na primeira subida
\end{minipage} \\
\midrule\noalign{}
\endhead
\bottomrule\noalign{}
\endlastfoot
\textbf{Patentes} & OPS (Escritório Europeu de Patentes), Lens Patent &
OPS \\
\textbf{Artigos} & Scopus, Lens Scholarly & Scopus \\
\end{longtable}
}

Também existe a \textbf{OpenAlex} (pública, sem chave), usada como
complemento para obter resumos que a Scopus não devolve para algumas
chaves. Ela não exige configuração.

\section{5.1 Passo a passo}\label{passo-a-passo}

\begin{enumerate}
\def\labelenumi{\arabic{enumi}.}
\tightlist
\item
  \textbf{Obtenha as credenciais} (são pessoais; não as envie a
  repositórios nem por canais inseguros):

  \begin{itemize}
  \tightlist
  \item
    \textbf{OPS:} consumer key e secret, no portal de desenvolvedores do
    Escritório Europeu de Patentes.
  \item
    \textbf{Scopus:} API key, no portal de desenvolvedores da Elsevier.
  \item
    \textbf{Lens:} token, em lens.org (serve para patentes e artigos).
  \end{itemize}
\item
  \textbf{Preencha} as credenciais no
  \texttt{.}\allowbreak{}\texttt{env} antes da primeira subida (elas são
  copiadas para o banco nessa subida), ou depois em Configurações ›
  Busca e credenciais:
  \texttt{OPS\_}\allowbreak{}\texttt{CONSUMER\_}\allowbreak{}\texttt{KEY},
  \texttt{OPS\_}\allowbreak{}\texttt{CONSUMER\_}\allowbreak{}\texttt{SECRET},
  \texttt{SCOPUS\_}\allowbreak{}\texttt{API\_}\allowbreak{}\texttt{KEY}
  e
  \texttt{LENS\_}\allowbreak{}\texttt{API\_}\allowbreak{}\texttt{TOKEN}.
\item
  \textbf{Escolha a API ativa} de cada família na tela Busca (ou via
  API, ver \hyperref[8-configurauxe7uxe3o-pela-api]{seção 8}).
\item
  \textbf{Reinicie o backend.} Esse passo é obrigatório: os adaptadores
  de busca são montados na subida. Sem o reinício, a API que você acabou
  de configurar \textbf{não aparece} como disponível, mesmo com a chave
  preenchida.
\end{enumerate}

\section{5.2 O que acontece sem
credencial}\label{o-que-acontece-sem-credencial}

\begin{itemize}
\tightlist
\item
  A família sem credencial é \textbf{pulada} na montagem (aviso
  \texttt{container\_}\allowbreak{}\texttt{*\_}\allowbreak{}\texttt{api\_}\allowbreak{}\texttt{skipped\ …\ missing\_}\allowbreak{}\texttt{credentials}
  no log).
\item
  O sistema sobe normalmente, mas a busca externa daquela família volta
  \textbf{vazia}.
\item
  Com nenhuma API de patentes e nenhuma de artigos configurada, nenhuma
  prospecção retorna resultados.
\end{itemize}

\section{5.3 Como uma API é
escolhida}\label{como-uma-api-uxe9-escolhida}

\begin{itemize}
\tightlist
\item
  Família \textbf{Patentes}: a seleção guarda a API ativa (\texttt{ops}
  ou \texttt{lens\_}\allowbreak{}\texttt{patent}).
\item
  Família \textbf{Artigos}: idem, com \texttt{scopus} ou
  \texttt{lens\_}\allowbreak{}\texttt{scholarly}.
\item
  A tela Busca mostra as opções e o estado de cada uma.
\end{itemize}

\begin{center}\rule{0.5\linewidth}{0.5pt}\end{center}

\chapter{6. Parâmetros da tela
Configurações}\label{paruxe2metros-da-tela-configurauxe7uxf5es}

A tela \textbf{Configurações} tem seis grupos. Cada parâmetro abaixo é
editável, gravado no banco e aplicado ao vivo (exceto os marcados como
reinício). Os padrões são os do código; o valor \textbf{vigente} é o que
aparece na tela.

\section{6.1 Geral}\label{geral}

{\def\LTcaptype{none} % do not increment counter
\begin{longtable}[]{@{}
  >{\raggedright\arraybackslash}p{(\linewidth - 6\tabcolsep) * \real{0.1667}}
  >{\raggedright\arraybackslash}p{(\linewidth - 6\tabcolsep) * \real{0.2000}}
  >{\raggedright\arraybackslash}p{(\linewidth - 6\tabcolsep) * \real{0.2500}}
  >{\raggedright\arraybackslash}p{(\linewidth - 6\tabcolsep) * \real{0.3833}}@{}}
\toprule\noalign{}
\begin{minipage}[b]{\linewidth}\raggedright
Parâmetro
\end{minipage} & \begin{minipage}[b]{\linewidth}\raggedright
Chave
\end{minipage} & \begin{minipage}[b]{\linewidth}\raggedright
Padrão
\end{minipage} & \begin{minipage}[b]{\linewidth}\raggedright
Para que serve
\end{minipage} \\
\midrule\noalign{}
\endhead
\bottomrule\noalign{}
\endlastfoot
Modelo de teste & \texttt{test\_}\allowbreak{}\texttt{mode} &
\texttt{false} & Usa respostas falsas em todas as chamadas de IA. Não
ligar em uso normal. \\
Modelo do KeyBERT &
\texttt{llm\_}\allowbreak{}\texttt{keybert\_}\allowbreak{}\texttt{model}
&
\texttt{distiluse-}\allowbreak{}\texttt{base-}\allowbreak{}\texttt{multilingual-}\allowbreak{}\texttt{cased-}\allowbreak{}\texttt{v2}
& Modelo de embeddings da extração de termos e do RAG. \textbf{Reinício
obrigatório.} \\
Timeout do Ollama (s) &
\texttt{ollama\_}\allowbreak{}\texttt{request\_}\allowbreak{}\texttt{timeout\_}\allowbreak{}\texttt{seconds}
& \texttt{600} & Espera máxima por chamadas ao servidor Ollama ou
compatível. \\
Corte relativo do RAG &
\texttt{rag\_}\allowbreak{}\texttt{relative\_}\allowbreak{}\texttt{min\_}\allowbreak{}\texttt{relevance}
& \texttt{0.}\allowbreak{}\texttt{75} & Descarta trechos recuperados
muito abaixo do melhor. \\
Top-K do RAG por seção &
\texttt{rag\_}\allowbreak{}\texttt{top\_}\allowbreak{}\texttt{k\_}\allowbreak{}\texttt{per\_}\allowbreak{}\texttt{section}
& \texttt{5} & Quantos trechos entram no prompt de cada seção. \\
Idioma da revisão de texto &
\texttt{languagetool\_}\allowbreak{}\texttt{language} &
\texttt{pt-}\allowbreak{}\texttt{BR} & Idioma do botão Revisão do
editor. \\
\end{longtable}
}

\section{6.2 Busca}\label{busca}

{\def\LTcaptype{none} % do not increment counter
\begin{longtable}[]{@{}
  >{\raggedright\arraybackslash}p{(\linewidth - 6\tabcolsep) * \real{0.1311}}
  >{\raggedright\arraybackslash}p{(\linewidth - 6\tabcolsep) * \real{0.4098}}
  >{\raggedright\arraybackslash}p{(\linewidth - 6\tabcolsep) * \real{0.0492}}
  >{\raggedright\arraybackslash}p{(\linewidth - 6\tabcolsep) * \real{0.4098}}@{}}
\toprule\noalign{}
\begin{minipage}[b]{\linewidth}\raggedright
Parâmetro
\end{minipage} & \begin{minipage}[b]{\linewidth}\raggedright
Chave
\end{minipage} & \begin{minipage}[b]{\linewidth}\raggedright
Padrão
\end{minipage} & \begin{minipage}[b]{\linewidth}\raggedright
Para que serve
\end{minipage} \\
\midrule\noalign{}
\endhead
\bottomrule\noalign{}
\endlastfoot
Ano inicial padrão &
\texttt{search\_}\allowbreak{}\texttt{year\_}\allowbreak{}\texttt{from}
& \texttt{2015} & Início da janela da busca exploratória. Quanto mais
antigo, mais lenta e menos focada. \\
Ano final padrão &
\texttt{search\_}\allowbreak{}\texttt{year\_}\allowbreak{}\texttt{to} &
\texttt{2026} & Fim da janela. Não deve ser maior que o ano atual. \\
Top-K da probe &
\texttt{probe\_}\allowbreak{}\texttt{top\_}\allowbreak{}\texttt{k} &
\texttt{10} & Resultados da busca exploratória. Aumentar encarece a
extração de termos. \\
Top-K da busca final &
\texttt{final\_}\allowbreak{}\texttt{top\_}\allowbreak{}\texttt{k} &
\texttt{100} & Resultados da busca final para os agregados. Muito alto
pode esbarrar em paginação das APIs. \\
Credenciais de busca &
\texttt{lens\_}\allowbreak{}\texttt{api\_}\allowbreak{}\texttt{token},
\texttt{ops\_}\allowbreak{}\texttt{consumer\_}\allowbreak{}\texttt{key},
\texttt{ops\_}\allowbreak{}\texttt{consumer\_}\allowbreak{}\texttt{secret},
\texttt{scopus\_}\allowbreak{}\texttt{api\_}\allowbreak{}\texttt{key} &
vazio & Ver \hyperref[5-apis-de-busca]{seção 5}. \\
\end{longtable}
}

\section{6.3 Inteligência Artificial}\label{inteliguxeancia-artificial}

Tela própria, com provedores e pontos de uso. Ver
\hyperref[4-provedores-de-ia]{seção 4}.

Os modelos iniciais dos provedores não são parâmetros da tela: são
gravados no cadastro do provedor na primeira subida e podem ser
alterados lá.

{\def\LTcaptype{none} % do not increment counter
\begin{longtable}[]{@{}
  >{\raggedright\arraybackslash}p{(\linewidth - 4\tabcolsep) * \real{0.1167}}
  >{\raggedright\arraybackslash}p{(\linewidth - 4\tabcolsep) * \real{0.2333}}
  >{\raggedright\arraybackslash}p{(\linewidth - 4\tabcolsep) * \real{0.6500}}@{}}
\toprule\noalign{}
\begin{minipage}[b]{\linewidth}\raggedright
Provedor
\end{minipage} & \begin{minipage}[b]{\linewidth}\raggedright
Modelo inicial
\end{minipage} & \begin{minipage}[b]{\linewidth}\raggedright
Observação
\end{minipage} \\
\midrule\noalign{}
\endhead
\bottomrule\noalign{}
\endlastfoot
\texttt{ollama} &
\texttt{qwen2.}\allowbreak{}\texttt{5:}\allowbreak{}\texttt{3b-}\allowbreak{}\texttt{instruct}
& Cadastrado sempre. O endereço vem de
\texttt{OLLAMA\_}\allowbreak{}\texttt{BASE\_}\allowbreak{}\texttt{URL}. \\
\texttt{gemini} &
\texttt{gemini-}\allowbreak{}\texttt{2.}\allowbreak{}\texttt{0-}\allowbreak{}\texttt{flash-}\allowbreak{}\texttt{exp}
& Cadastrado só se
\texttt{LLM\_}\allowbreak{}\texttt{GEMINI\_}\allowbreak{}\texttt{API\_}\allowbreak{}\texttt{KEY}
estiver preenchida. \\
\texttt{anthropic} &
\texttt{claude-}\allowbreak{}\texttt{haiku-}\allowbreak{}\texttt{4-}\allowbreak{}\texttt{5}
& Cadastrado só se
\texttt{LLM\_}\allowbreak{}\texttt{ANTHROPIC\_}\allowbreak{}\texttt{API\_}\allowbreak{}\texttt{KEY}
estiver preenchida. \\
\end{longtable}
}

\section{6.4 Relevância \& Qualidade}\label{relevuxe2ncia-qualidade}

{\def\LTcaptype{none} % do not increment counter
\begin{longtable}[]{@{}
  >{\raggedright\arraybackslash}p{(\linewidth - 6\tabcolsep) * \real{0.2500}}
  >{\raggedright\arraybackslash}p{(\linewidth - 6\tabcolsep) * \real{0.2333}}
  >{\raggedright\arraybackslash}p{(\linewidth - 6\tabcolsep) * \real{0.0500}}
  >{\raggedright\arraybackslash}p{(\linewidth - 6\tabcolsep) * \real{0.4667}}@{}}
\toprule\noalign{}
\begin{minipage}[b]{\linewidth}\raggedright
Parâmetro
\end{minipage} & \begin{minipage}[b]{\linewidth}\raggedright
Chave
\end{minipage} & \begin{minipage}[b]{\linewidth}\raggedright
Padrão
\end{minipage} & \begin{minipage}[b]{\linewidth}\raggedright
Para que serve
\end{minipage} \\
\midrule\noalign{}
\endhead
\bottomrule\noalign{}
\endlastfoot
Threshold de relevância &
\texttt{relevance\_}\allowbreak{}\texttt{threshold} &
\texttt{0.}\allowbreak{}\texttt{4} & Corte mínimo para um resultado ser
considerado válido. \\
Complexidade máxima de query &
\texttt{llm\_}\allowbreak{}\texttt{max\_}\allowbreak{}\texttt{query\_}\allowbreak{}\texttt{complexity}
& \texttt{0.}\allowbreak{}\texttt{6} & Teto de complexidade de query
gerada por IA. \\
\end{longtable}
}

\section{6.5 Extração de Termos}\label{extrauxe7uxe3o-de-termos}

{\def\LTcaptype{none} % do not increment counter
\begin{longtable}[]{@{}
  >{\raggedright\arraybackslash}p{(\linewidth - 6\tabcolsep) * \real{0.2034}}
  >{\raggedright\arraybackslash}p{(\linewidth - 6\tabcolsep) * \real{0.2881}}
  >{\raggedright\arraybackslash}p{(\linewidth - 6\tabcolsep) * \real{0.0508}}
  >{\raggedright\arraybackslash}p{(\linewidth - 6\tabcolsep) * \real{0.4576}}@{}}
\toprule\noalign{}
\begin{minipage}[b]{\linewidth}\raggedright
Parâmetro
\end{minipage} & \begin{minipage}[b]{\linewidth}\raggedright
Chave
\end{minipage} & \begin{minipage}[b]{\linewidth}\raggedright
Padrão
\end{minipage} & \begin{minipage}[b]{\linewidth}\raggedright
Efeito de aumentar
\end{minipage} \\
\midrule\noalign{}
\endhead
\bottomrule\noalign{}
\endlastfoot
Peso do título &
\texttt{term\_}\allowbreak{}\texttt{extraction\_}\allowbreak{}\texttt{title\_}\allowbreak{}\texttt{weight}
& \texttt{3.}\allowbreak{}\texttt{0} & Termos dependem mais do
título. \\
Peso do abstract &
\texttt{term\_}\allowbreak{}\texttt{extraction\_}\allowbreak{}\texttt{abstract\_}\allowbreak{}\texttt{weight}
& \texttt{1.}\allowbreak{}\texttt{0} & Abstract domina sobre o
título. \\
Threshold de score &
\texttt{term\_}\allowbreak{}\texttt{extraction\_}\allowbreak{}\texttt{score\_}\allowbreak{}\texttt{threshold}
& \texttt{0.}\allowbreak{}\texttt{024} & Menos termos, só os melhores.
Se o lote inteiro ficar abaixo, entra o piso abaixo. \\
Piso de termos devolvidos &
\texttt{term\_}\allowbreak{}\texttt{extraction\_}\allowbreak{}\texttt{min\_}\allowbreak{}\texttt{returned\_}\allowbreak{}\texttt{terms}
& \texttt{10} & Mais termos garantidos, possivelmente de qualidade
menor. \\
Penalidade de unigrama &
\texttt{term\_}\allowbreak{}\texttt{extraction\_}\allowbreak{}\texttt{unigram\_}\allowbreak{}\texttt{penalty}
& \texttt{-}\allowbreak{}\texttt{0.}\allowbreak{}\texttt{4} & Mais
negativo favorece termos compostos. \\
Bônus de bigrama &
\texttt{term\_}\allowbreak{}\texttt{extraction\_}\allowbreak{}\texttt{bigram\_}\allowbreak{}\texttt{bonus}
& \texttt{0.}\allowbreak{}\texttt{0} & Favorece bigramas quando alto. \\
Bônus de trigrama &
\texttt{term\_}\allowbreak{}\texttt{extraction\_}\allowbreak{}\texttt{trigram\_}\allowbreak{}\texttt{bonus}
& \texttt{0.}\allowbreak{}\texttt{25} & Favorece frases mais longas e
específicas. \\
Penalidade de bigrama ruim &
\texttt{term\_}\allowbreak{}\texttt{extraction\_}\allowbreak{}\texttt{bad\_}\allowbreak{}\texttt{bigram\_}\allowbreak{}\texttt{penalty}
& \texttt{-}\allowbreak{}\texttt{0.}\allowbreak{}\texttt{8} & Filtra
mais padrões gramaticais ruins. \\
Penalidade de trigrama ruim &
\texttt{term\_}\allowbreak{}\texttt{extraction\_}\allowbreak{}\texttt{bad\_}\allowbreak{}\texttt{trigram\_}\allowbreak{}\texttt{penalty}
& \texttt{-}\allowbreak{}\texttt{0.}\allowbreak{}\texttt{8} & Idem, para
trigramas. \\
BM25 k1 &
\texttt{term\_}\allowbreak{}\texttt{extraction\_}\allowbreak{}\texttt{bm25\_}\allowbreak{}\texttt{k1}
& \texttt{1.}\allowbreak{}\texttt{2} & Termos muito repetidos dominam
mais. \\
BM25 b &
\texttt{term\_}\allowbreak{}\texttt{extraction\_}\allowbreak{}\texttt{bm25\_}\allowbreak{}\texttt{b}
& \texttt{0.}\allowbreak{}\texttt{75} & Penaliza mais documentos
longos. \\
RRF k &
\texttt{term\_}\allowbreak{}\texttt{extraction\_}\allowbreak{}\texttt{rrf\_}\allowbreak{}\texttt{k}
& \texttt{60} & Valores altos suavizam a fusão de rankings. \\
Frequência mínima (C-value) &
\texttt{term\_}\allowbreak{}\texttt{extraction\_}\allowbreak{}\texttt{cvalue\_}\allowbreak{}\texttt{min\_}\allowbreak{}\texttt{frequency}
& \texttt{1} & Descarta termos raros. \\
\end{longtable}
}

\section{6.6 Inferência Estatística}\label{inferuxeancia-estatuxedstica}

{\def\LTcaptype{none} % do not increment counter
\begin{longtable}[]{@{}
  >{\raggedright\arraybackslash}p{(\linewidth - 6\tabcolsep) * \real{0.2131}}
  >{\raggedright\arraybackslash}p{(\linewidth - 6\tabcolsep) * \real{0.3443}}
  >{\raggedright\arraybackslash}p{(\linewidth - 6\tabcolsep) * \real{0.0492}}
  >{\raggedright\arraybackslash}p{(\linewidth - 6\tabcolsep) * \real{0.3934}}@{}}
\toprule\noalign{}
\begin{minipage}[b]{\linewidth}\raggedright
Parâmetro
\end{minipage} & \begin{minipage}[b]{\linewidth}\raggedright
Chave
\end{minipage} & \begin{minipage}[b]{\linewidth}\raggedright
Padrão
\end{minipage} & \begin{minipage}[b]{\linewidth}\raggedright
Para que serve
\end{minipage} \\
\midrule\noalign{}
\endhead
\bottomrule\noalign{}
\endlastfoot
Duração máxima (s) &
\texttt{statistical\_}\allowbreak{}\texttt{inference\_}\allowbreak{}\texttt{max\_}\allowbreak{}\texttt{duration\_}\allowbreak{}\texttt{seconds}
& \texttt{60} & Tempo para pedir buscas extras até a amostra saturar. \\
Threshold de saturação &
\texttt{statistical\_}\allowbreak{}\texttt{inference\_}\allowbreak{}\texttt{saturation\_}\allowbreak{}\texttt{threshold}
& \texttt{0.}\allowbreak{}\texttt{5} & Cobertura mínima estimada. \\
Threshold de singletons &
\texttt{statistical\_}\allowbreak{}\texttt{inference\_}\allowbreak{}\texttt{f1\_}\allowbreak{}\texttt{ratio\_}\allowbreak{}\texttt{threshold}
& \texttt{0.}\allowbreak{}\texttt{7} & Proporção máxima de itens vistos
uma vez. \\
Mínimo de doubletons &
\texttt{statistical\_}\allowbreak{}\texttt{inference\_}\allowbreak{}\texttt{f2\_}\allowbreak{}\texttt{min}
& \texttt{5} & Mínimo de itens vistos duas vezes. \\
Reamostragens (bootstrap) &
\texttt{statistical\_}\allowbreak{}\texttt{inference\_}\allowbreak{}\texttt{bootstrap\_}\allowbreak{}\texttt{resamples}
& \texttt{1000} & Precisão da estabilidade do ranking. Mais = mais
lento. \\
Máximo de títulos (relevância) &
\texttt{statistical\_}\allowbreak{}\texttt{inference\_}\allowbreak{}\texttt{max\_}\allowbreak{}\texttt{titles\_}\allowbreak{}\texttt{for\_}\allowbreak{}\texttt{relevance}
& \texttt{20} & Títulos usados na relevância semântica. \\
\end{longtable}
}

\section{6.7 Fuzzy matching}\label{fuzzy-matching}

{\def\LTcaptype{none} % do not increment counter
\begin{longtable}[]{@{}
  >{\raggedright\arraybackslash}p{(\linewidth - 6\tabcolsep) * \real{0.2000}}
  >{\raggedright\arraybackslash}p{(\linewidth - 6\tabcolsep) * \real{0.2667}}
  >{\raggedright\arraybackslash}p{(\linewidth - 6\tabcolsep) * \real{0.0500}}
  >{\raggedright\arraybackslash}p{(\linewidth - 6\tabcolsep) * \real{0.4833}}@{}}
\toprule\noalign{}
\begin{minipage}[b]{\linewidth}\raggedright
Parâmetro
\end{minipage} & \begin{minipage}[b]{\linewidth}\raggedright
Chave
\end{minipage} & \begin{minipage}[b]{\linewidth}\raggedright
Padrão
\end{minipage} & \begin{minipage}[b]{\linewidth}\raggedright
Para que serve
\end{minipage} \\
\midrule\noalign{}
\endhead
\bottomrule\noalign{}
\endlastfoot
Threshold de fuzzy match &
\texttt{depositant\_}\allowbreak{}\texttt{fuzzy\_}\allowbreak{}\texttt{match\_}\allowbreak{}\texttt{threshold}
& \texttt{90.}\allowbreak{}\texttt{0} & Similaridade para tratar nomes
parecidos como a mesma entidade. Mais baixo agrupa mais (risco de juntar
entidades diferentes). \\
\end{longtable}
}

\begin{center}\rule{0.5\linewidth}{0.5pt}\end{center}

\chapter{7. Relatório: capa, LaTeX e
revisão}\label{relatuxf3rio-capa-latex-e-revisuxe3o}

\section{7.1 Imagem de capa}\label{imagem-de-capa}

Tela \textbf{Configurações › Geral › Imagem de capa do relatório}, com
recorte. Use \textbf{Trocar imagem} para enviar outra capa, ou
\textbf{Restaurar imagem padrão} para voltar ao brasão do sistema. A
capa padrão é enviada ao MinIO automaticamente na subida, se ainda não
houver capa. Uma capa enviada pelo usuário nunca é sobrescrita pelo
sistema. A imagem recortada é usada na capa do REPTEC.

API: \texttt{GET}, \texttt{POST} e
\texttt{DELETE\ /}\allowbreak{}\texttt{config/}\allowbreak{}\texttt{report-}\allowbreak{}\texttt{cover-}\allowbreak{}\texttt{image}
(o \texttt{DELETE} restaura a capa padrão).

\section{7.2 Compilação de PDF}\label{compilauxe7uxe3o-de-pdf}

O PDF é gerado pelo serviço
\texttt{latex-}\allowbreak{}\texttt{compiler} (container), sob demanda.
Limites do compilador:

\begin{itemize}
\tightlist
\item
  Roda \textbf{sem internet}: só usa os pacotes instalados no serviço.
\item
  Usa \textbf{pdflatex}
  (\texttt{latexmk\ -}\allowbreak{}\texttt{pdf\ -}\allowbreak{}\texttt{halt-}\allowbreak{}\texttt{on-}\allowbreak{}\texttt{error},
  limite de 120 s por compilação).
\item
  Não suporta \texttt{fontspec} nem XeLaTeX/LuaLaTeX; caracteres fora do
  alfabeto latino não compilam.
\end{itemize}

\textbf{Pacotes disponíveis:} \texttt{graphicx}, \texttt{float},
\texttt{tabularx}, \texttt{array}, \texttt{enumitem}, \texttt{caption},
\texttt{hyperref}, \texttt{setspace} e \texttt{babel} (português), além
dos que o próprio template carrega (\texttt{geometry},
\texttt{fancyhdr}, \texttt{microtype}, \texttt{lastpage},
\texttt{titlesec}, \texttt{tocloft}, \texttt{indentfirst},
\texttt{tgtermes}). As coleções LaTeX recomendadas e de fontes também
estão instaladas.

\textbf{Pacotes ausentes:} \texttt{multirow}, \texttt{makecell},
\texttt{siunitx}, \texttt{wrapfig}, \texttt{ulem} e \texttt{soul}. Um
\texttt{\textbackslash{}usepackage} desses interrompe a compilação com
``File \ldots{} .sty not found''.

Para adicionar um pacote ao compilador, inclua-o na lista de instalação
do serviço \texttt{latex-}\allowbreak{}\texttt{compiler} e reconstrua
esse serviço. Isso exige internet no momento do build.

\section{7.3 Revisão de texto}\label{revisuxe3o-de-texto}

O botão \textbf{Revisão} usa o LanguageTool
(\texttt{LANGUAGETOOL\_}\allowbreak{}\texttt{URL}, idioma
\texttt{languagetool\_}\allowbreak{}\texttt{language}) e a IA para
sugerir correções. Cada sugestão é aplicada ou ignorada individualmente.
O sistema recusa sugestões que alterem tempo verbal, pontuação ou
maiúsculas.

\begin{center}\rule{0.5\linewidth}{0.5pt}\end{center}

\chapter{8. Configuração pela API}\label{configurauxe7uxe3o-pela-api}

Todas as rotas ficam sob
\texttt{/}\allowbreak{}\texttt{api/}\allowbreak{}\texttt{v1}. A
documentação interativa fica em \url{http://localhost:8000/docs}.

{\def\LTcaptype{none} % do not increment counter
\begin{longtable}[]{@{}
  >{\raggedright\arraybackslash}p{(\linewidth - 2\tabcolsep) * \real{0.4500}}
  >{\raggedright\arraybackslash}p{(\linewidth - 2\tabcolsep) * \real{0.5500}}@{}}
\toprule\noalign{}
\begin{minipage}[b]{\linewidth}\raggedright
Método e rota
\end{minipage} & \begin{minipage}[b]{\linewidth}\raggedright
Função
\end{minipage} \\
\midrule\noalign{}
\endhead
\bottomrule\noalign{}
\endlastfoot
\texttt{GET\ /}\allowbreak{}\texttt{config/}\allowbreak{}\texttt{settings}
& Lista todas as configurações editáveis, com valor, tipo, categoria e
limites. \\
\texttt{PUT\ /}\allowbreak{}\texttt{config/}\allowbreak{}\texttt{settings/}\allowbreak{}\texttt{\{key\}}
& Altera uma configuração. Corpo:
\texttt{\{"value":}\allowbreak{}\texttt{\ ".}\allowbreak{}\texttt{.}\allowbreak{}\texttt{.}\allowbreak{}\texttt{"\}}.
\textbf{O valor é sempre texto}, mesmo para números e booleanos (ex.:
\texttt{"2015"}, \texttt{"false"}). Aplicada imediatamente. \\
\texttt{GET\ /}\allowbreak{}\texttt{config/}\allowbreak{}\texttt{search-}\allowbreak{}\texttt{apis}
& Mostra as opções de API por família e a ativa. \\
\texttt{PUT\ /}\allowbreak{}\texttt{config/}\allowbreak{}\texttt{search-}\allowbreak{}\texttt{apis/}\allowbreak{}\texttt{\{family\}}
& Troca a API ativa. Corpo:
\texttt{\{"active\_}\allowbreak{}\texttt{provider\_}\allowbreak{}\texttt{code":}\allowbreak{}\texttt{\ "ops"\}}.
\textbf{Reinício obrigatório.} \\
\texttt{GET\ /}\allowbreak{}\texttt{config/}\allowbreak{}\texttt{llm/}\allowbreak{}\texttt{providers}
& Provedores suportados e se exigem chave ou URL. \\
\texttt{GET\ /}\allowbreak{}\texttt{config/}\allowbreak{}\texttt{llm/}\allowbreak{}\texttt{configs}
& Configurações de IA cadastradas (chave mascarada). \\
\texttt{POST\ /}\allowbreak{}\texttt{config/}\allowbreak{}\texttt{llm/}\allowbreak{}\texttt{configs}
& Cadastra uma configuração de IA. \\
\texttt{PUT\ /}\allowbreak{}\texttt{config/}\allowbreak{}\texttt{llm/}\allowbreak{}\texttt{configs/}\allowbreak{}\texttt{\{id\}}
& Altera uma configuração de IA. \\
\texttt{DELETE\ /}\allowbreak{}\texttt{config/}\allowbreak{}\texttt{llm/}\allowbreak{}\texttt{configs/}\allowbreak{}\texttt{\{id\}}
& Remove uma configuração de IA. \\
\texttt{GET\ /}\allowbreak{}\texttt{config/}\allowbreak{}\texttt{llm/}\allowbreak{}\texttt{call-}\allowbreak{}\texttt{sites}
& Mostra qual configuração atende cada ponto de uso. \\
\texttt{PUT\ /}\allowbreak{}\texttt{config/}\allowbreak{}\texttt{llm/}\allowbreak{}\texttt{call-}\allowbreak{}\texttt{sites/}\allowbreak{}\texttt{\{call\_}\allowbreak{}\texttt{site\}}
& Aponta um ponto de uso para outra configuração. \\
\texttt{GET}, \texttt{POST},
\texttt{DELETE\ /}\allowbreak{}\texttt{config/}\allowbreak{}\texttt{report-}\allowbreak{}\texttt{cover-}\allowbreak{}\texttt{image}
& Imagem de capa do relatório. \\
\end{longtable}
}

Exemplo --- alterar o ano inicial da busca:

\begin{Shaded}
\begin{Highlighting}[]
\FunctionTok{curl} \OperatorTok{{-}}\NormalTok{X PUT http}\OperatorTok{://}\NormalTok{localhost}\OperatorTok{:}\DecValTok{8000}\OperatorTok{/}\NormalTok{api}\OperatorTok{/}\NormalTok{v1}\OperatorTok{/}\NormalTok{config}\OperatorTok{/}\NormalTok{settings}\OperatorTok{/}\NormalTok{search\_year\_from \textasciigrave{}}
  \OperatorTok{{-}}\NormalTok{H }\StringTok{"Content{-}Type: application/json"}\NormalTok{ \textasciigrave{}}
  \OperatorTok{{-}}\NormalTok{d }\StringTok{\textquotesingle{}\{"value": "2015"\}\textquotesingle{}}
\end{Highlighting}
\end{Shaded}

Exemplo --- trocar a API de patentes para OPS:

\begin{Shaded}
\begin{Highlighting}[]
\FunctionTok{curl} \OperatorTok{{-}}\NormalTok{X PUT http}\OperatorTok{://}\NormalTok{localhost}\OperatorTok{:}\DecValTok{8000}\OperatorTok{/}\NormalTok{api}\OperatorTok{/}\NormalTok{v1}\OperatorTok{/}\NormalTok{config}\OperatorTok{/}\NormalTok{search{-}apis}\OperatorTok{/}\NormalTok{patent \textasciigrave{}}
  \OperatorTok{{-}}\NormalTok{H }\StringTok{"Content{-}Type: application/json"}\NormalTok{ \textasciigrave{}}
  \OperatorTok{{-}}\NormalTok{d }\StringTok{\textquotesingle{}\{"active\_provider\_code": "ops"\}\textquotesingle{}}
\end{Highlighting}
\end{Shaded}

\begin{center}\rule{0.5\linewidth}{0.5pt}\end{center}

\chapter{9. Segurança}\label{seguranuxe7a}

Itens que \textbf{você deve tratar antes de uso compartilhado ou em
intranet}:

\begin{enumerate}
\def\labelenumi{\arabic{enumi}.}
\tightlist
\item
  \textbf{Senhas padrão.} Os serviços do Docker vêm com senhas de
  desenvolvimento: PostgreSQL (\texttt{postgres}/\texttt{postgres}),
  MinIO (\texttt{minioadmin}/\texttt{minioadmin}) e pgAdmin
  (\texttt{admin@example.}\allowbreak{}\texttt{com}/\texttt{admin}).
  Troque antes de expor qualquer porta.
\item
  \textbf{Chaves no banco.} Chaves de API de IA e de busca ficam
  gravadas no PostgreSQL. Restrinja o acesso ao banco e não faça backup
  em local aberto.
\item
  \textbf{Arquivo \texttt{.}\allowbreak{}\texttt{env} fora do
  repositório.} O \texttt{.}\allowbreak{}\texttt{env} real contém
  chaves. Não o envie a repositórios; confira antes de cada envio de
  código.
\item
  \textbf{\texttt{SECRET\_}\allowbreak{}\texttt{KEY} e
  \texttt{ALGORITHM} não são usados.} O sistema não tem autenticação por
  token, então esses valores não têm efeito. Não os trate como proteção.
\item
  \textbf{Modo de teste desligado.}
  \texttt{TEST\_}\allowbreak{}\texttt{MODE=}\allowbreak{}\texttt{False}
  em uso real. Com ele ligado, a aplicação responde com dados falsos e
  não chama nenhuma IA.
\item
  \textbf{CORS.} Em \texttt{ALLOWED\_}\allowbreak{}\texttt{ORIGINS},
  deixe apenas as origens que de fato acessam a API.
\item
  \textbf{Exposição de portas.} Os serviços do Docker publicam portas no
  host. Em uma máquina compartilhada, limite a interface de cada porta
  (ex.:
  \texttt{127.}\allowbreak{}\texttt{0.}\allowbreak{}\texttt{0.}\allowbreak{}\texttt{1:}\allowbreak{}\texttt{5432:}\allowbreak{}\texttt{5432})
  na configuração dos serviços.
\end{enumerate}

\begin{center}\rule{0.5\linewidth}{0.5pt}\end{center}

\chapter{10. Quando reiniciar o
backend}\label{quando-reiniciar-o-backend}

A maior parte das alterações vale na hora. \textbf{Reinicie o backend}
nestes casos:

{\def\LTcaptype{none} % do not increment counter
\begin{longtable}[]{@{}
  >{\raggedright\arraybackslash}p{(\linewidth - 2\tabcolsep) * \real{0.4667}}
  >{\raggedright\arraybackslash}p{(\linewidth - 2\tabcolsep) * \real{0.5333}}@{}}
\toprule\noalign{}
\begin{minipage}[b]{\linewidth}\raggedright
Alteração
\end{minipage} & \begin{minipage}[b]{\linewidth}\raggedright
Motivo
\end{minipage} \\
\midrule\noalign{}
\endhead
\bottomrule\noalign{}
\endlastfoot
Credenciais de OPS, Scopus ou Lens (novas ou trocadas) & Os adaptadores
de busca são montados na subida. \\
Troca da API ativa de patentes ou artigos & Idem. O backend registra no
log o evento
\texttt{search\_}\allowbreak{}\texttt{api\_}\allowbreak{}\texttt{selection\_}\allowbreak{}\texttt{updated\_}\allowbreak{}\texttt{restart\_}\allowbreak{}\texttt{required}. \\
Modelo do KeyBERT
(\texttt{llm\_}\allowbreak{}\texttt{keybert\_}\allowbreak{}\texttt{model})
& Carregado uma vez no boot. \\
Alteração no \texttt{.}\allowbreak{}\texttt{env} (credenciais e flags) &
Lido na subida. Parâmetros de comportamento não são lidos do
\texttt{.}\allowbreak{}\texttt{env} (ver
\hyperref[2-ciclo-de-vida-do-env-ao-banco]{seção 2}). \\
\texttt{ALLOWED\_}\allowbreak{}\texttt{ORIGINS},
\texttt{DATABASE\_}\allowbreak{}\texttt{URL},
\texttt{API\_}\allowbreak{}\texttt{PREFIX} & Configuração de subida. \\
\end{longtable}
}

Para reiniciar: interrompa o processo do backend (\texttt{Ctrl+C}) e
suba de novo.

\begin{center}\rule{0.5\linewidth}{0.5pt}\end{center}

\chapter{11. Exemplos de cenários}\label{exemplos-de-cenuxe1rios}

Os ajustes abaixo são feitos na tela \textbf{Configurações} (não no
\texttt{.}\allowbreak{}\texttt{env}).

\textbf{Cenário A --- Máquina sem acesso a IA externa (padrão).}
Provedor \texttt{ollama} para as queries, com
\texttt{OLLAMA\_}\allowbreak{}\texttt{BASE\_}\allowbreak{}\texttt{URL=}\allowbreak{}\texttt{http:}\allowbreak{}\texttt{/}\allowbreak{}\texttt{/}\allowbreak{}\texttt{localhost:}\allowbreak{}\texttt{11434}
e o modelo baixado. Não precisa de chave de IA. Precisa de chave de
busca para obter resultados.

\textbf{Cenário B --- Usar Gemini para as queries e Ollama para o texto
do relatório.} 1. Cadastre um provedor \texttt{gemini} com a chave em
Configurações › Inteligência Artificial. 2. Em \textbf{Qual IA cada
chamada usa}, aponte \texttt{theme\_}\allowbreak{}\texttt{candidates},
\texttt{probe\_}\allowbreak{}\texttt{query} e
\texttt{final\_}\allowbreak{}\texttt{query} para o Gemini. 3. Deixe
\texttt{report\_}\allowbreak{}\texttt{writing} e
\texttt{report\_}\allowbreak{}\texttt{review} no Ollama.

\textbf{Cenário C --- Servidor de IA da intranet.} 1. Cadastre um
provedor \texttt{ollama} com a Base URL do servidor e, se houver, a API
key. 2. Aponte os pontos de uso para essa configuração. 3. Ajuste o
timeout do Ollama se o servidor for lento.

\textbf{Cenário D --- Busca mais focada (menos ruído).} - Aumente o
threshold de relevância (ex.: 0,5) e o corte relativo do RAG. - Reduza o
ano inicial se o tema for recente.

\textbf{Cenário E --- Busca mais ampla (mais resultados, mais ruído).} -
Reduza o threshold de relevância. - Aumente o Top-K da busca final
(atenção ao tempo e aos limites das APIs).

\textbf{Cenário F --- Termos mais curtos e genéricos ou mais longos e
específicos.} - Mais longos: aumente o bônus de trigrama e deixe a
penalidade de unigrama mais negativa. - Mais curtos: reduza o bônus de
trigrama e aproxime a penalidade de unigrama de zero.

\begin{center}\rule{0.5\linewidth}{0.5pt}\end{center}

\chapter{12. Limitações conhecidas}\label{limitauxe7uxf5es-conhecidas}

\begin{itemize}
\tightlist
\item
  \textbf{Valores do \texttt{.}\allowbreak{}\texttt{env} após a primeira
  subida.} Editar no \texttt{.}\allowbreak{}\texttt{env} parâmetros de
  comportamento não tem efeito, e editar credenciais também não: ambos
  já foram gravados no banco. Use a tela.
\item
  \textbf{Credenciais novas exigem reinício.} Ver
  \hyperref[10-quando-reiniciar-o-backend]{seção 10}.
\item
  \textbf{Status do Ollama.} A consulta de provedor atual pode indicar
  \texttt{available:}\allowbreak{}\texttt{\ true} mesmo com o Ollama
  parado. Teste com
  \texttt{curl\ http:}\allowbreak{}\texttt{/}\allowbreak{}\texttt{/}\allowbreak{}\texttt{localhost:}\allowbreak{}\texttt{11434/}\allowbreak{}\texttt{api/}\allowbreak{}\texttt{tags}.
\item
  \textbf{Rodapé da barra lateral.} Em versões anteriores, mostrava
  valores fixos (``LLM: GEMINI 2.5'', ``API Artigos: SCOPUS'') e não
  refletia a configuração atual. Para conferir qual IA está em uso,
  consulte a tela Configurações.
\item
  \textbf{Primeira subida precisa de internet} para baixar o modelo de
  embeddings.
\item
  \textbf{Modelo do spaCy.} Ele não é instalado pelos pacotes Python do
  projeto; instale-o conforme o passo de ambiente Python descrito na
  instalação. Sem ele, a extração de termos funciona, mas com resultados
  diferentes.
\end{itemize}

\end{document}
